\subsubsection{The FKG inequality for Gaussian vectors}

The FKG inequality is a crucial tool to apply percolation arguments. We say that a Borel subset $A\subseteq\R^n$ is \emph{increasing} if for each $x \in A$ and $y\in\R^n$ such that $y_i \geq x_i$ for every $i \in \lbrace 1,\dots, n \rbrace$, we have $y \in A$. We say that $A$ is decreasing if the complement of $A$ is increasing. The FKG inequality for Gaussian vectors was proved by Pitt and can be stated as follows.

\begin{theo}[\cite{pitt1982positively}]\label{t.FKG}
Let $X = (X_1,\dots, X_n)$ be a $n$-dimensional Gaussian vector whose correlation matrix has non-negative entries. Then, for every increasing Borel subsets $A,B \subseteq \R^n$, we have:
\[
\Pro \left[X \in A \cap B \right] \geq \Pro \left[X \in A \right] \cdot \Pro \left[X \in B \right] \,.
\]
\end{theo}

This result is the reason why we work with positively correlated Gaussian fields. Indeed, the FKG inequality is a crucial ingredient in the proof of RSW-type results. Note however that the very recent~\cite{bg_17} proves a box-crossing property without this inequality, albeit only in a discrete setting.

\subsubsection{A differential formula}\label{sss.russo}

As in the case of Bernoulli percolation, we need to introduce a notion of influence. This notion of influence is inspired by the \emph{geometric influences} studied by Keller, Mossel and Sen in the case of product spaces~\cite{keller2012geometric,keller2014geometric}. (For more about the relations between the geometric influences and the influences of Definition~\ref{d.infl} below, which are roughly the same in the case of product measures, see Subsection~\ref{ss.KMSnonproduct}.)

\begin{defi}\label{d.infl}
Let $\mu$ be a finite Borel measure on $\R^n$, let $v \in \R^n$, and let $A$ be a Borel subset of $\R^n$. The influence of $v$ on $A$ under $\mu$ is:
\[
I_{v,\mu}(A) := \underset{r \downarrow 0}{\liminf} \, \frac{\mu \left(A + [-r,r] \, v \right) - \mu \left(A \right)}{r} \in [0,+\infty] \,.
\]
Write $\left(e_1,\dots, e_n \right)$ for the canonical basis of $\R^n$. We will use the following simplified notations:
\[
I_{i,\mu}(A) := I_{e_i,\mu}(A) \,.
\]
\end{defi}
